\chapter{Uygulanan Test Yaklaşımı}

\section{Data-Driven Test Tasarımı (JSON Senaryolar)}
Her modül için JSON dosyasında senaryo tanımları bulunur. Örnek alanlar: \texttt{caseId}, \texttt{description}, \texttt{technique}, \texttt{standard}, \texttt{testType}, \texttt{expectedOutcome}.

\begin{lstlisting}[caption={Örnek senaryo alanları (login-scenarios.json)}]
{
  "caseId": "LGN-001",
  "description": "Gecersiz sifre ile giris",
  "technique": "negative-test",
  "standard": "OWASP",
  "testType": "regression",
  "testLevel": "system-test"
}
\end{lstlisting}

\section{Risk Bazlı Senaryo Hedefleri}
Modül karmaşıklığı ve iş kritikliğine göre hedef senaryo sayıları belirlenmiştir:

\begin{table}[H]
  \centering
  \caption{Risk bazlı modül hedefleri}
  \begin{tabular}{@{}lr@{}}
    \toprule
    \textbf{Modül} & \textbf{Hedef} \\
    \midrule
    Checkout & 28 \\
    Login / ForgotPassword & 24 \\
    Register / Cart / MyAccountSettings & 22 \\
    MyAccount & 20 \\
    Search & 18 \\
    WishList & 12 \\
    Newsletter & 10 \\
    Api & 32 \\
    Db (H2 + SQL Server) & 35 \\
    E2EOrderFlow & 4 \\
    \midrule
    \textbf{Toplam} & \textbf{$\sim$273} \\
    \bottomrule
  \end{tabular}
\end{table}

\section{Test Teknikleri}
Equivalence Partitioning (EP), Boundary Value Analysis (BVA), negatif/pozitif test, durum geçişi ve güvenlik odaklı senaryolar uygulanmıştır.

\section{OWASP ve ISTQB Standart Eşlemesi}
Dashboard Modül Kapsamı bölümünde OWASP test kategorileri modüllere eşlenmiştir. Senaryo metadata'sında \texttt{standard} alanı ile referans tutulur.

\section{UI Modül Kapsamı}
10 UI modülü: Login, Register, ForgotPassword, MyAccount, MyAccountSettings, Newsletter, WishList, Search, Cart, Checkout. Toplam $\sim$202 UI senaryosu.

\section{API, DB ve E2E Senaryo Kapsamı}
\begin{itemize}
  \item \textbf{API:} GET/POST/PUT/DELETE, auth, pagination, hata kodları (32 senaryo).
  \item \textbf{DB:} CRUD, constraint, transaction, SQL Server entegrasyon (35 senaryo).
  \item \textbf{E2E:} UI checkout $\rightarrow$ API $\rightarrow$ DB doğrulama (4 senaryo).
\end{itemize}

% TODO: Modül başına pass/fail tablosu ekle (son koşum sonuçlarından).
